
\section{Motivation}
\label{sec:motivation}

TODO: explain \lstinline{drop}

\subsection{Borrowing from Context-Free Session Types}
\label{sec:borr-from-cont}


Context-free session types enable writing elegant code for, e.g.,
serializing tree structures.%
\footnote{We assume session-typed communication over a network
  or between sandboxed processes, where only primitive data can be
  passed, but no pointers.}
Consider an example from \citet{DBLP:conf/icfp/ThiemannV16} for serializing binary trees of type \lstinline{Tree}.
\begin{lstlisting}
type Tree = Leaf
          | Node int Tree Tree
\end{lstlisting}
Here is the corresponding session type \lstinline{TreeChannel}:
\begin{lstlisting}
type TreeChannel = oplus{
  Leaf: skip,
  Node: !int; TreeChannel; TreeChannel
}
\end{lstlisting}
The type \lstinline{TreeChannel} is called context-free because it is
recursively defined and the recursion is not restricted to tail
positions. The \lstinline{Node} alternative is an example where the
\lstinline{TreeChannel} type refers two times in a row to itself. This
type would be illegal in regular session type systems (e.g.,
\cite{DBLP:journals/jfp/GayV10,DBLP:conf/esop/LindleyM15}) because
they syntactically restrict recursion to tail position. Context-free
session types lift this restriction as their composition operator $S_1;S_2$ composes
two session types sequentially: the composed session comprises the
actions of $S_1$ first, then followed by the actions of $S_2$.\footnote{One
  might argue that CFST are close to Honda's original vision for
  sesssion types \cite{DBLP:conf/concur/Honda93}, which also features
  symmetric composition.}

Given these types, the serialization function looks straightforward:
\begin{lstlisting}[numbers=left]
sendTree : Tree -> forallalpha.TreeChannel;alpha -> alpha
sendTree Leaf c =
  select Leaf c
sendTree (Node x l r) c =
  let c1 = select Node c   (* c1 : !int; TreeChannel; TreeChannel; alpha *)
      c2 = send x c1       (* c2 : TreeChannel; TreeChannel; alpha *)
      c3 = sendTree l c2   (* c3 : TreeChannel; alpha *)
      c4 = sendTree r c3   (* c4 : alpha *)
  in  c4
\end{lstlisting}
However, there are several catches. First, the channel \lstinline{c}
is passed explicitly throughout the code; a pattern we call
\emph{resource-passing style}. This style forces the programmer to
make the function polymorphic on the continuation type of the channel,
which is \lstinline{alpha} in this example.
Second, the \lstinline{Node} branch of \lstinline{sendTree} relies on
polymorphic recursion. To see this, consider the type of
\lstinline{c2} in line 6--7. To invoke \lstinline{sendTree}
recursively, we need to instantiate the type argument
\lstinline{alpha} with \lstinline{TreeChannel; alpha} rather than just
\lstinline{alpha}, which is the
telltale for polymorphic recursion. Finally, recall that channel types
are linear in CFST, that is, each session-typed variable must be used
exactly once.

Our new language {\thecalculus} amends both problems. The type
definitions remain unchanged, but the implementation of
\lstinline{sendTree} changes.
\begin{lstlisting}[numbers=left]
sendTree' : Tree -> TreeChannel; Ret -> Unit
sendTree' Leaf c =          (* c : TreeChannel; Ret *)
  drop (select Leaf c)      (* ... : Ret *)
sendTree' (Node x l r) c =
  let c = select Node c in  (* c : !Int; TreeChannel; TreeChannel; Ret *)
  send x &c;                (* c : TreeChannel; TreeChannel; Ret *)
  sendTree l &c;            (* c : TreeChannel; Ret *)
  sendTree r c
\end{lstlisting}
What has changed? First, the type is no longer polymorphic over the
continuation type. Instead, the channel argument \lstinline{c} starts
out at type \lstinline{TreeChannel; Ret} which indicates that it is a borrow from
a (larger) session type. Second, there is no explicit resource
passing, except for \lstinline{select}. Each use of the borrow notation
\lstinline{&c} implicitly splits off an initial segment of the session type
according to the function argument and updates \lstinline{c} to the
remainder. To see this, consider the type of \lstinline{c}---after
splitting off the borrow in the corresponding line---shown in the
comments. Before the \lstinline{send} operation (line 5--6), the type
of \lstinline{c} starts with \lstinline{!Int; ...}. ``Taking the
borrow'' splits off the borrowed type \lstinline{!Int; Ret}---the
argument type of \lstinline{send x}---and
updates \lstinline{c} to the type shown in line 6. As the type of
\lstinline{c} is itself a borrowed type, one could speak of a
\emph{reborrow}. On the last line~8, the type of \lstinline{c} already
matches the type expected by \lstinline{sendTree}, so no extra borrow
needs to be taken. 

Third, as there is no polymorphism, there is no polymorphic
recursion, either.

Finally, the channel variable \lstinline{c} can be used in a
non-linear way, as long as its uses are under borrows, as in
\lstinline{&c}. The understanding is that the borrows are taken and
processed according to the evaluation order of the program. 

The reader may ask why \lstinline{select} still adheres to
resource passing style. The answer is simplicity of our system. While it
is possible to define \lstinline{select} with a borrowing-based
interface, supporting it requires subtyping of session types, which is
undecidable for CFST. 

\subsection{Local and Remote Borrowing}
\label{sec:local-remote-borr}

We call the style of borrowing demonstrated in
\cref{sec:borr-from-cont} \emph{local borrowing}, as the borrows remain in the same
thread and are processed sequentially, following the evaluation order.

In this local setting, we expect the operations on borrows to be
implemented efficiently in an imperative style. That is, the channel is
represented by a reference to some data structure (e.g., a file
descriptor) and each channel operation updates 
the data structure. This implementation of local borrows is safe because the soundness of
the type system ensures that the type updates match the real updates
in the implementation.

However, there may be situations where we want to send a borrow to
another process. Suppose, for example, a server-side renderer streams
HTML nodes to a client instead of constructing a complete DOM tree in
memory first. A small protocol for such a renderer is as follows.
\begin{lstlisting}
type HtmlChannel = oplus{
  Text: !String,
  Elem: !String; ChildrenChannel
}
type ChildrenChannel = oplus{
  Done: skip,
  Child: HtmlChannel; ChildrenChannel
}
\end{lstlisting}
The \lstinline{Elem} alternative first sends a tag name, and then a sequence of children. The 
\lstinline{Child} alternative is context-free in the same sense as the
\lstinline{Node} alternative of \lstinline{TreeChannel}: rendering one
child consumes an \lstinline{HtmlChannel} prefix and then continues with
the rest of the child stream.

Now consider a renderer that emits a user panel whose first child is
computed by another process. We assume that
\lstinline{renderProfile} and \lstinline{renderSwiper} both render
one HTML node on a channel of type \lstinline{HtmlChannel; End}.
\Cref{listing:rendering-example} contains the code.
\begin{lstlisting}[numbers=left,float,caption={Rendering example},captionpos=b,label=listing:rendering-example]
renderUser : User -> HtmlChannel; End -> Unit
renderUser user c =
  let c = select Elem c in
  send "section" &c;
  let c = select Child c in
  fork (lambda() -> renderProfile user &c);
  renderUserRest user (acquire c)

renderUserRest : User -> ChildrenChannel; End -> Unit
renderUserRest user c =
  let c = select Child c in
  renderSwiper user &c;
  close (select Done c)
\end{lstlisting}
The function \lstinline{renderUser} starts an HTML element and then
opens the child stream. It passes a borrow of the
first child position to a new process, which renders the user's profile.
The continuation renders the remaining children, here a permissions
view, and then closes the child stream.

In this case, the borrow under the \lstinline{fork} is a \emph{remote
  borrow}. Such a borrow \emph{cannot} be handled imperatively: the
continuation must not start emitting the second child before the
forked renderer has finished the first child. This synchronization is
implemented with the \lstinline{acquire} function. It blocks until the
borrow created by the fork is completed, i.e., until
\lstinline{renderProfile} has finished sending on the borrowed prefix
of the channel. Once it is unblocked, the \lstinline{acquire}
operation returns the current channel end for further operations.
Next, we a direct operational semantics of local and remote borrowing
before we discuss the low level implementation of the two concepts in
\cref{sec:impl-local-remote}.

\subsection{A Direct Semantics for Channel Borrowing}
\label{sec:direct-semant-chann}

\Cref{sec:borr-from-cont} already suggests an implementation strategy
for local borrowing. However, we would like to start with an abstract
specification that is easier to handle in a correctness proof, before moving on to
the implementation level. The first step translates the surface
language with borrowing operators into our formal calculus
\thecalculus, which replaces borrowing with explicit functional splitting
operators. This translation is taken from previous work
\cite{DBLP:journals/pacmpl/SaffrichS0V25}, which gives a formal
presentation along with a correctness proof. The novelty in our work
is the presence of \emph{two} splitting operators, one for local borrows and
one for remote ones. Our discussion starts with the local variant and
we show the translation of \lstinline{renderUserRest} in
\thecalculus. The translation is part of our implementation.

As an example, we use the local part of
\lstinline{renderUserRest} as shown in
\cref{fig:local-split-program}.  
Execution proceeds by repeatedly exposing the next local prefix of the
channel sequence. \Cref{fig:local-split-trace} shows a suffix of its reduction
sequence. To keep the trace short, the
reduction starts after the initial local splits and the selection of
\lstinline{Child}; the omitted steps only expose the prefix consumed by
\lstinline{renderSwiper}.
\begin{figure}
\begin{subfigure}[t]{0.49\linewidth}
\begin{lstlisting}[numbers=left]
K(c) =
  let c2 = select Child c in
  let c3, c4 = lsplit c2 in
  renderSwiper user c3;
  close (select Done c4)
\end{lstlisting}
\caption{Translated program.}
\label{fig:local-split-program}
\end{subfigure}\hfill
\begin{subfigure}[t]{0.49\linewidth}
\begin{lstlisting}
nu[c][d..] in
  (P_d ||
   let c2 = select Child c in
   let c3, c4 = lsplit c2 in
   renderSwiper used c3;
   close (select Done c4)
~> nu[c2][d..] in
  (P_d ||
   let c3, c4 = lsplit c2 in
   renderSwiper user c3;
   close (select Done c4)
~> nu[c3,c4][d..] in
  (P_d ||
   renderSwiper user c3;
   close (select Done c4)
~> nu[c4][d..] in
  (P_d ||
   close (select Done c4)
~> nu[c4][d..] in
  (P_d ||
   close c4
~> nu[c4][d3] in
  (P_d'[wait d3] ||
   close c4)
~> (P_d'[*] || *)
\end{lstlisting}
\caption{Reduction suffix.}
\label{fig:local-split-trace}
\end{subfigure}
\caption{Local splitting in the suffix of \lstinline{renderUserRest}.
The processes $P_d$, $P_d'$, and $P_d''$ abbreviate client-side
computation on the other endpoint.}
\label{fig:local-split-reduction}
\end{figure}
In the omitted prefix, the split on \lstinline{c2} \emph{replaces} the
binder for \lstinline{c2} with two fresh channel names \lstinline{c3}
and \lstinline{c4}. That is, both binding compartments of the
\lstinline{nu} contain a \emph{list} of binders and the operations
work on these channel names from left to right. To ensure the correct
sequencing, the semantics does not allow any action on \lstinline{c4} before
\lstinline{c3} has been consumed. This restriction is supported by
typing: the \lstinline{lsplit} operation returns an \emph{ordered
  pair}, whose components must be processed and consumed in
left-to-right order.

On the client side, in the process $P_d$ that holds the other
channel end \lstinline{d}, 
splitting can be performed on \lstinline{d} independently of
the splitting on \lstinline{c}. The only constraint is that, in the
end, when \lstinline{close} and \lstinline{wait} calls match, the
reduction requires that only one channel name is left in each
compartment. Under this condition, the channel can be closed and
eliminated. The figure does not show the transitions of the client $P_d$.

In principle, this semantics would also work for remote
borrowing. However, to connect to the \lstinline{drop}
\lstinline{acquire} primitives, we indicate a remote split with the
symbol ``\lstinline{|}''. As an example, we consider the translation of
lines 5--7:
\begin{lstlisting}
  fork (lambda() -> renderProfile user &c);
  renderUserRest user (acquire c)
\end{lstlisting}
It translates to the following \thecalculus code.
\begin{lstlisting}
(* initially                  c : HtmlChannel; ChildrenChannel; End *)
let c1, c2 = rsplit c in   (* c1 : HtmlChannel; Ret *)
                           (* c2 : Acquire; ChildrenChannel; End *)
fork (lambda() -> renderProfile user c1);
renderUserRest user (acquire c2)
\end{lstlisting}
The translation inserts an \lstinline{rsplit} because \lstinline{c1}
and \lstinline{c2} are used in different processes. As their uses are
explicitly synchronized, the \lstinline{rsplit} returns a normal
(linear) pair of channels. The type of the second channel
\lstinline{c2} starts with an \lstinline{Acquire} type, a special form
of session type that indicates the requirement for explicit
synchronization.
Executing this code takes the following steps (ignoring the activity
at the other end of the channel).
%     nu[c][...]
%     let c = select Child c in
%     let c1, c2 = rsplit c in
%     fork (lambda() -> renderProfile user c1);
%     renderUserRest user (acquire c2)
% ~>

\begin{lstlisting}
    nu[c][...]
    let c1, c2 = rsplit c in
    fork (lambda() -> renderProfile user c1);
    renderUserRest user (acquire c2)
~>
    nu[c1|c2][...]
    fork (lambda() -> renderProfile user c1);
    renderUserRest user (acquire c2)
~>
    nu[|c2][...]
    renderUserRest user (acquire c2)
~>
    nu[c2][...]
    renderUserRest user c2
\end{lstlisting}
In summary, the \lstinline{rsplit} inserts a new binding group \lstinline{[c1|c2]}
demarkated by ``\lstinline{|}''. The \lstinline{acquire} is only
enables if the preceding binding group has been processed as in
\lstinline{[|c2]}. Executing the \lstinline{acquire} removes the
exhausted binding group as in \lstinline{[c2]}.

\subsection{Implementing Local and Remote Borrowing}
\label{sec:impl-local-remote}


%%% Local Variables:
%%% mode: latex
%%% TeX-master: "../popl2027.tex"
%%% End:
